\section{实验}

\subsection{数据与评测}
本文在 PersonPath22 \cite{shuai2022personpath22} 上研究重连，使用其原生人工关键帧，以及发布方在 CrowdHuman 上训练、按帧与标注来源对齐的 FCOS 行人检测结果。观测网格名义上为 5~Hz。236 条公开片段中，16 条发布方托管的训练资产划为开发集，91 条划为校准集，另留 74 条发布方托管的官方测试资产用于评测；55 条带外部来源映射的片段不进入本协议。所有跟踪器变体使用相同的检测输入、匹配规则与评测量序列。

设定刻意限定在只有检测的证据上。不训练重识别网络，也不保存特征库：门控只能依据几何、检测置信度和分配裕量作判断。外观重识别，例如 BoT-SORT 的重识别变体 \cite{aharon2022botsort}，携带的是正交信号，而且往往更强；它是本文的自然延伸，不是本文的组成部分。

为度量检测缺口恢复，先把检测框与真值行人按 IoU $\geq 0.50$ 一对一匹配。对每个真值身份，缺口定义为同一身份两个相邻匹配检测之间有 1 到 5 个连续标注关键帧、且这些帧均无匹配检测。若缺口后第一个回归检测的轨迹编号与缺口前一致，计为一次恢复。所有策略共用同一份缺口清单，缺口前未输出过身份的样本也计入分母。

重连动作与缺口恢复分开审计。若指定的一对一匹配把前一个已输出轨迹与新检测框归到同一标注行人，该次重连正确；否则计为错误或无法核验，但未匹配端点与忽略区域重叠超过自身框面积一半的情形除外，这些动作单列为忽略。审计在跟踪之后进行，门控本身看不到任何标注。整体指标为 HOTA、AssA、IDF1 与编号切换数，在原生关键帧网格上用 TrackEval \cite{luiten2021hota} 计算，并采用 PersonPath22 的干扰物与人群处理。

\subsection{对照与拟合}
回放 ByteTrack \cite{zhang2022bytetrack} 与仅用运动的 BoT-SORT \cite{aharon2022botsort}，均从第一帧开始，门控在跟踪器更新状态之前决定每一次原生旧身份重连。对照对象包括不限额跟踪器、仅限额、原生关联代价阈值、最近竞争者裕量阈值、加性线性评分、交互评分，以及带同源排序项的交互评分。限额策略每来源最多接受 12 次动作，忽略动作同样占用限额。评分在开发集上做 200 次 Adam 更新（学习率 0.04，系数惩罚 0.002），显卡为 RTX 3060 Ti。候选阈值取自校准前就固定下来的开发集分数量点。每个评分取"在共同校准天花板内恢复缺口最多、且至少 20\% 校准来源活跃"的工作点，随后不做调整地运行于官方测试资产。

\subsection{官方划分结果}
官方测试子集含 9,320 张关键帧与 3,155 个可用检测缺口。在 47.23\% 的共同校准天花板下，ByteTrack 的交互门控把已测的接错或无法核验动作比例从 46.77\% 降到 35.39\%，核验正确的重连由 239 次升到 283 次，可评估接受数由 449 降到 438。恢复缺口分别为 393 与 388。按来源配对重采样的区间为 $[-16.57,-6.40]$ 个百分点（失败率差）与 $[-0.33,0.67]$ 个百分点（缺口恢复差），10,000 次重采样。在同一天花板下，交互评分相对仅限额与线性对照也都降低了失败比例。

在 20\% 天花板下，两种 ByteTrack 评分都恢复 383 个缺口，实测失败比例为 21.24\%（交互）与 21.15\%（线性）。运动版 BoT-SORT 上，交互评分恢复 334 个缺口，线性为 327，核验失败比例分别为 16.19\% 与 18.31\%，配对区间为 $[0.05,0.48]$（缺口恢复）与 $[-4.94,0.75]$（失败比例）。在 42.15\% 的松弛天花板上，线性评分的实测失败比例反而更低，为 25.19\% 对 37.28\%。加入同源排序项没有带来一致收益。

\begin{figure*}[t]
\centering
\begin{tikzpicture}
\begin{axis}[
    width=0.86\textwidth,
    height=0.40\textwidth,
    xlabel={恢复的自然缺口（共 3,155 个）},
    ylabel={实测失败比例（\%）},
    xmin=315, xmax=545,
    ymin=12, ymax=64,
    grid=major,
    grid style={gray!18},
    legend pos=north east,
    legend cell align=left,
    legend style={draw=none, fill=none, font=\small},
    tick label style={font=\small},
    label style={font=\small},
    mark size=2.4pt,
]
\addplot[only marks, mark=o, mark options={solid}, color=blue!65!black,
         forget plot] coordinates {
    (529,59.95) (399,57.73) (388,46.77) (391,44.30) (395,48.05)
    (393,35.39) (393,47.15) (383,21.15) (383,21.24)
};
\addlegendimage{only marks, mark=o, color=blue!65!black}
\addlegendentry{ByteTrack}
\addplot[only marks, mark=square, mark options={solid}, color=red!65!black,
         forget plot] coordinates {
    (424,53.39) (354,48.69) (343,37.76) (347,41.02) (339,25.19)
    (346,37.28) (344,37.78) (327,18.31) (334,16.19)
};
\addlegendimage{only marks, mark=square, color=red!65!black}
\addlegendentry{BoT-SORT（仅运动）}
\addplot[only marks, mark=star, mark size=5pt, color=black] coordinates {(529,59.95)};
\addplot[only marks, mark=star, mark size=5pt, color=black] coordinates {(424,53.39)};
\addplot[only marks, mark=triangle*, mark size=4pt, color=blue!65!black] coordinates {(393,35.39)};
\addplot[only marks, mark=triangle*, mark size=4pt, color=red!65!black] coordinates {(334,16.19)};
\node[font=\scriptsize, anchor=south west] at (axis cs:531,59.0) {不限额};
\node[font=\scriptsize, anchor=north east] at (axis cs:391,34.2) {交互};
\node[font=\scriptsize, anchor=south east] at (axis cs:332,16.9) {严格};
\end{axis}
\end{tikzpicture}
\caption{官方 PersonPath22 测试子集上的全部已评估策略，横轴为恢复的自然缺口数，
纵轴为可评估重连中的实测失败比例。星号标出不限额跟踪器，三角标出松弛交互门控，
左下角一组为严格天花板。靠右的策略恢复更多缺口，靠下的策略更少失败。前沿从左下的
严格工作点延伸到右上的不限额原版，没有任何一个策略在两个轴上都最优。}
\label{fig:tradeoff-zh}
\end{figure*}


\begin{table*}[t]
\centering
\small
\caption{PersonPath22 原生关键帧官方测试子集。所有策略使用同一组 3,155 个缺口。接受数与失败数已剔除忽略动作。HOTA、AssA、IDF1 为 0--100 分制，IDSW 为计数，风险列为实测失败比例。仅限额门控使用同样的来源上限。未标"严格"的行使用松弛校准天花板。第~\ref{sec:matched-zh} 节在匹配接受量规则下重新选点，其计数与本表不同。}
\label{tab:reconnection-main-zh}
\begin{tabular}{llrrrrrrrr}
\hline
跟踪器 & 重连策略 & 保留缺口 & 接受 & 失败 & 风险 (\%) & HOTA & AssA & IDF1 & IDSW \\
\hline
OC-SORT & 外部参照 & 174 & -- & -- & -- & 41.83 & 41.22 & 52.15 & 3790 \\
\hline
ByteTrack & 不限额 & 529 & 2589 & 1552 & 59.95 & 46.64 & 42.33 & 56.73 & 6662 \\
 & 仅限额 & 399 & 537 & 310 & 57.73 & 44.69 & 39.32 & 52.94 & 7129 \\
 & 原生代价 & 388 & 449 & 210 & 46.77 & 44.76 & 39.41 & 53.16 & 7104 \\
 & 最近裕量 & 391 & 456 & 202 & 44.30 & 44.84 & 39.57 & 53.14 & 7132 \\
 & 加性 & 395 & 514 & 247 & 48.05 & 44.90 & 39.63 & 53.38 & 7131 \\
 & 交互 & 393 & 438 & 155 & 35.39 & 44.99 & 39.81 & 53.38 & 7022 \\
 & 交互 + 排序 & 393 & 509 & 240 & 47.15 & 44.91 & 39.69 & 53.34 & 7131 \\
 & 加性（严格） & 383 & 331 & 70 & 21.15 & 44.87 & 39.63 & 53.25 & 7120 \\
 & 交互（严格） & 383 & 339 & 72 & 21.24 & 44.87 & 39.63 & 53.23 & 7116 \\
\hline
BoT-SORT & 不限额 & 424 & 1858 & 992 & 53.39 & 46.84 & 43.78 & 56.57 & 5909 \\
（仅运动） & 仅限额 & 354 & 495 & 241 & 48.69 & 45.04 & 40.71 & 53.03 & 6616 \\
 & 原生代价 & 343 & 392 & 148 & 37.76 & 44.99 & 40.67 & 52.98 & 6656 \\
 & 最近裕量 & 347 & 451 & 185 & 41.02 & 45.10 & 40.83 & 53.07 & 6654 \\
 & 加性 & 339 & 397 & 100 & 25.19 & 45.20 & 40.99 & 53.29 & 6603 \\
 & 交互 & 346 & 448 & 167 & 37.28 & 45.14 & 40.88 & 53.20 & 6566 \\
 & 交互 + 排序 & 344 & 450 & 170 & 37.78 & 45.14 & 40.87 & 53.16 & 6611 \\
 & 加性（严格） & 327 & 284 & 52 & 18.31 & 45.05 & 40.83 & 53.05 & 6628 \\
 & 交互（严格） & 334 & 278 & 45 & 16.19 & 45.06 & 40.85 & 53.02 & 6626 \\
\hline
\end{tabular}
\end{table*}

不限额跟踪器保留了更多自然缺口，整体 HOTA 也更高（表~\ref{tab:reconnection-main-zh}）。相对于不限额的 ByteTrack，松弛交互门控减少了 1,397 次可评估失败，也少保留 136 个自然缺口，HOTA 下降 1.65 分。我们把这个差额当作选择性接受的代价报出来，而不是设法把它解释掉：门控买到的是决策可靠性，不是一个更好的跟踪器。OC-SORT \cite{cao2023ocsort} 作为整体跟踪参照列入，其输出支持同一套缺口评测，但它的内部重连机制走的是另一套协议。

\subsection{接受量对齐}
\label{sec:matched-zh}

只数错误会奖励沉默。一个什么都不接受的门控没有失败，而表~\ref{tab:reconnection-main-zh} 比较的策略接受量从 278 到 2,589 次不等，排序里因此混入了接受量的影响。我们直接去掉这个混淆：对每个跟踪器和每个上限 $B\in\{6,12,24\}$，重新选阈值，使原生关联代价、加性评分与交互评分在每一来源上接受同样多的动作。

对齐是有代价的。目标是校准集上每来源的接受数，允许的偏离为每来源 0.10 次或目标值的 5\%，取较大者。原始分位网格的 18 个组合没有一个达标，于是把扫描加密到 0.01 步长并在校准集上回放，直到能锁定 17 个工作点。其中 7 个在测试来源上仍满足方法间容差，其余 10 个作为敏感性结果报告。17 个里面每一个都至少有一个方法偏离校准目标超过容差，所以这是方法之间的比较，不是对某个固定接受率的复现。

对原生关联代价，交互评分在全部 17 个工作点上失败比例更低，差距为 8.70 到 23.00 个百分点，每个比较各自的区间都不含零。在上限 12、高档这一具体设置上：ByteTrack 是 381 次可评估接受里失败 98 次，原生代价为 390 次里失败 158 次；仅用运动的 BoT-SORT 是 367 次里失败 81 次，原生代价为 350 次里失败 116 次；配对区间分别为 $[-20.95,-8.80]$ 与 $[-17.29,-5.15]$ 个百分点。缺口恢复的变化为 $+0.41$ $[+0.03,+0.84]$ 与 $+0.10$ $[-0.35,+0.60]$ 个百分点。也就是说，可靠性提升在这个工作点上没有以恢复为代价，它也不是"少接受"造成的假象。

对加性评分，情况要薄一些。17 个对照里只有两个区间完全低于零：ByteTrack 高档的 $-3.29$ $[-5.37,-1.10]$ 与 BoT-SORT 在 $B=24$ 时的 $-2.57$ $[-5.05,-0.16]$；其余 15 个跨零，其中若干点估计为正（ByteTrack 中档 $+1.54$ $[-1.83,+4.85]$）。对齐接受量之后，交互评分与原生代价的分野很干净，与加性模型的分野只出现在部分工作点。我们保留这一区别，而不是把它平均掉。

\subsection{增益来源}
\label{sec:ablation-zh}

有两个问题决定评分是贡献还是拟合产物：$\phi$ 的哪一部分承载了效果，以及已发表的关联评分是否能用更简单的办法做到同样的事。

对第一个问题，我们在同一批开发提案、标签和样本顺序上构造三个评分。12 维与完整的 23 维评分沿用已锁定的权重；15 维评分保留 12 个基础特征再加三个竞争摘要（两个裕量的最小值、裕量之差的绝对值、变换后备选计数的最大值），并且从头重新拟合，因为截断完整模型的权重无法区分这两个假设。拟合沿用原设置，包括显式把排序权重设为零。我们在同一提案上核对了 15 维特征前缀与完整向量前 15 项逐项相等，并确认重新载入的 12 维与 23 维配置复现各自存档的分数。

表~\ref{tab:round3-zh} 给出配对对照。把 8 个乘积项加到 15 维评分上，ByteTrack 高档的失败比例下降 3.29 个百分点，区间 $[-5.82,-0.92]$；中档方向相反，为 $+1.60$ $[-1.43,+4.48]$；BoT-SORT 三档区间全部跨零。三个竞争摘要单独使用时，在五个工作点上没有任何一个对照区间排除零。效果归属于这组乘积项整体，并且在单个工作点得到确证，而不是贯穿全部工作点。

\begin{table*}[t]
\centering
\small
\caption{交互项消融与外部评分对照，每来源上限 12 次。每行是同一批 74 个来源上的一个配对对照，重采样 10,000 次（种子 362）。“失败/接受”为剔除忽略区域后的可评估决策数，不是缺口恢复的分母。$\Delta$ 为左策略失败比例减右策略失败比例（百分点），负值有利于左。BoostTrack 行是在同一宿主内移植关联评分后的比较。}
\label{tab:round3-zh}
\begin{tabular}{lllrrr}
\hline
跟踪器 & 天花板 & 对照 & 失败/接受（左 / 右） & $\Delta$ (pp) & 95\% 区间 \\
\hline
ByteTrack & 高 & 23 对 15 & 98/381 \;/\; 112/386 & $-3.29$ & $[-5.82,-0.92]$ \\
 & 中 & 23 对 15 & 69/327 \;/\; 62/318 & $+1.60$ & $[-1.43,+4.48]$ \\
 & 高 & 23 对 BoostTrack & 98/381 \;/\; 152/388 & $-13.45$ & $[-19.06,-8.06]$ \\
 & 中 & 23 对 BoostTrack & 69/327 \;/\; 109/315 & $-13.50$ & $[-18.88,-8.40]$ \\
\hline
BoT-SORT & 低 & 23 对 15 & 30/222 \;/\; 33/225 & $-1.15$ & $[-3.02,+0.46]$ \\
（仅运动） & 中 & 23 对 15 & 57/293 \;/\; 54/291 & $+0.90$ & $[-1.56,+3.54]$ \\
 & 高 & 23 对 15 & 81/367 \;/\; 81/359 & $-0.49$ & $[-2.46,+1.40]$ \\
 & 低 & 23 对 BoostTrack & 30/222 \;/\; 53/214 & $-11.25$ & $[-17.32,-5.32]$ \\
 & 中 & 23 对 BoostTrack & 57/293 \;/\; 83/273 & $-10.95$ & $[-18.09,-3.84]$ \\
 & 高 & 23 对 BoostTrack & 81/367 \;/\; 110/346 & $-9.72$ & $[-16.00,-3.66]$ \\
\hline
\end{tabular}
\end{table*}

对第二个问题，我们把已发表跟踪器 BoostTrack \cite{stanjoevic2024boosttrack} 的关联相似度移植到同一角色。宿主跟踪器照原样分配；对它已经提出的候选，从当前帧候选与宿主轨迹状态读出被提议配对的 BoostTrack 相似度，再施加同一门控与同样的拒绝后处理。检测框、分数、原生指派、运动更新与后续阶段都不改动。评分组件被替换时容易与输入变化混淆，所以先做数值核对：在 1,747 个开发集候选矩阵（745,096 个元素）上，移植实现与作者函数的最大绝对误差为零；对仅用运动的 BoT-SORT——其状态保存的是宽与高而非高与宽高比——以协方差一阶近似传播后，420,620 个元素同样误差为零。

五个对照里，交互评分都比移植评分更少失败：ByteTrack 为 13.45 与 13.50 个百分点，BoT-SORT 为 9.72 到 11.25 个百分点，所有区间均低于零（表~\ref{tab:round3-zh}）。缺口恢复的差异不超过 0.25 个百分点，且这些区间全部跨零。在 ByteTrack 高档，交互评分、15 维评分与 BoostTrack 评分的 HOTA 分别为 44.98、44.97 与 44.78。这里的比较对象是共同宿主内的评分组件，整套 BoostTrack 跟踪器是另一个问题。整张表有一条限定：没有任何一个对照在测试时两端同时回到校准目标，BoT-SORT 低档的 15 对 12 还超出了方法间容差。阈值没有据此重选，报告的是锁定值。

\subsection{跨域迁移}
把 PersonPath22 的 ByteTrack 阈值迁移到 DanceTrack 训练划分下 40 条此前未见过的片段，使用冻结的 COCO YOLOX-X 检测。松弛交互门控恢复了 8,472 个缺口中的 1,506 个，462 次接受里失败 199 次（43.07\%）；原生代价门控恢复 1,500 个，463 次里失败 212 次（45.79\%）。在 20\% 校准工作点上，交互评分在 DanceTrack 的实测失败比例为 30.03\%。BDD100K 的五条探索视频中，严格工作点下的 ByteTrack 交互门控在 8 次接受里失败 3 次，不限额原版为 25 次里失败 16 次，两者都只保留 121 个可用检测缺口中的 5 个。在新的检测器和新的采集条件下，绝对失败比例明显高于 PersonPath22，但评分之间的排序延续了下来。

\subsection{代价与敏感性}
当匹配到的真值身份唯一时，我们对每个坐标在 0.001 像素以内数值等价的重复检测框做了重分配。这项事后检查没有改变松弛 ByteTrack 原生代价与交互对照、严格 BoT-SORT 加性与交互对照中的任何接受或失败计数。松弛 ByteTrack 交互门控在 \texttt{uid\_vid\_00013} 的关键帧 47--49 之间保留了标注身份 66，原生阈值则漏掉了这次回归；反例出现在 \texttt{uid\_vid\_00008}，身份 20 的关键帧 66--68 缺口之后，原生阈值接回了旧编号，而交互评分漏接。

关联与决策记录的中位耗时为每关键帧 1.60 与 1.54 毫秒（不限额与松弛交互 ByteTrack），以及 1.48 与 1.49 毫秒（不限额与严格交互 BoT-SORT）。计时覆盖七次 CPU 回放，每次 293 张开发关键帧，不含检测器推理与评测。门控对每个提案只是一次标量求值，开销不随轨迹数增长。

\subsection{适用范围}
表~\ref{tab:reconnection-main-zh} 之后的每一项分析都复用同一批 74 个资产。接受量对齐工作点、消融与外部评分对照都是同一组来源上的敏感性分析，不是彼此独立的重复实验，任何区间都没有为其余分析留出余量。门控不使用外观信号，这使结论限定在只有检测器的流水线上。每来源 12 次的上限是一个设计参数：在我们锁定的各档上限上方法排序保持，绝对数值不保持。跨域结果迁移的只是冻结阈值。
